% Supporting proof fragment; definitions and author--date bibliography are in the main paper.
\section{Certified published conditioning bounds and the classical information test}
\label{supp:classical-price-bounds}

I distinguish a continuous basket-price certificate from a continuous-to-grid
transfer certificate. The former encloses the unknown price directly. The
latter only bounds one error term after choosing a finite numerical law.
The published geometric exercise-set lower bound and arithmetic--geometric
upper bound are due to \citet{caldana}; I consulted the repository-labelled
accepted manuscript, Sections 1--3, Proposition 1.1, Proposition 2.1,
Corollary 2.2, equation (27) and Appendices B and D, rather than the
inaccessible final publisher body. The following complete derivation fixes
my Gaussian implementation, its mixed-sign optimiser and its capped-call
extension. It does not claim a new generic pricing theorem.

\subsection{The unchanged financial law and exact one-dimensional quantities}

Write $L_i=\log(S_i/K)$, $g=\sum_iw_iL_i$, $A=\sum_iw_iS_i$ and
$G=K e^g$, with $w_i>0$ and $\sum_iw_i=1$. Under the stipulated
risk-neutral Gaussian law, set
\[
 m_G=\mathbb E g,\quad v_G=\operatorname{Var}(g)>0,\quad
 \sigma_G=\sqrt{v_G},\quad
 a_i=\frac{\operatorname{Cov}(L_i,g)}{\sigma_G},\quad
 A_i=w_i\mathbb E S_i>0.
\]
All covariances are computed from the original rational exposures, including
both their common Gaussian loading and geometric outer loading. Positive
basket weights do not imply $a_i\geq0$. In particular, some $a_i$ are
negative in the retained two- and four-asset small-common-scale models.

The standardised variable $W=(g-m_G)/\sigma_G$ is standard normal.
The Gaussian variable $L_i-\mathbb E L_i-a_iW$ has covariance zero with
$W$. Their joint Gaussian characteristic function factorises, proving
independence, including a possibly degenerate residual variance.
It follows that
\[
 \mathbb E[S_i\mid W=z]
 =\mathbb E S_i\exp(a_i z-a_i^2/2).
\]
Indeed, the conditional residual variance is
$\operatorname{Var}(L_i)-a_i^2\geq0$ by Cauchy--Schwarz; evaluating its
Gaussian exponential moment gives exactly the displayed factor. Multiplying
by the standard normal density and completing the square yields
\[
 \mathbb E[S_i\mathbf1_{W>z}]=\mathbb E S_i\Phi(a_i-z),
 \qquad \mathbb P(W>z)=\Phi(-z).
\]
These quantities require no arithmetic-basket price as input.

For a strike $k>0$, define
\[
 L_k(z)=\sum_iA_i\Phi(a_i-z)-k\Phi(-z),\qquad
 \ell_k=\max\{0,\sup_{z\in\mathbb R}L_k(z)\}.
\]
Pointwise, $(A-k)\mathbf1_E\leq(A-k)_+$ for every event $E$.
Taking $E=\{W>z\}$ and expectations proves
$\ell_k\leq\mathbb E(A-k)_+$. The supremum includes the limiting
values $L_k(-\infty)=\mathbb E A-k$ and $L_k(+\infty)=0$.
Both limits follow by dominated convergence since $\mathbb E A<\infty$.

Jensen gives $A\geq G$. The call function is nondecreasing and
one-Lipschitz, so
\[
 (G-k)_+\leq(A-k)_+\leq(G-k)_++A-G.
\]
Consequently a valid upper bound is
$u_k=C_G(k)+\mathbb E A-\mathbb E G$, where direct Gaussian integration
gives
\[
 C_G(k)=K e^{m_G+v_G/2}
 \Phi\!\left(\frac{m_G-\log(k/K)+v_G}{\sigma_G}\right)
 -k\Phi\!\left(\frac{m_G-\log(k/K)}{\sigma_G}\right),
 \qquad \mathbb E G=K e^{m_G+v_G/2}.
\]
The same derivation gives the geometric lower bound $C_G(k)$.
Thus $[\ell_k,u_k]$ encloses the continuous undiscounted call.
The geometric lower bound is also valid, but no additional intersection
with it is credited to the frozen implemented comparator. A lower interval endpoint must be rounded downwards;
an upper endpoint must be rounded upwards. An interval containing the
lower bound is not itself a price interval until this direction is imposed.

\subsection{A complete optimiser, including mixed covariance signs}

Since $\phi(a_i-z)=\phi(z)e^{a_i z-a_i^2/2}$,
\[
 L_k'(z)=\phi(z)\{k-S(z)\},\quad
 S(z)=\sum_i A_i e^{a_i z-a_i^2/2},\quad
 S''(z)=\sum_i A_i a_i^2 e^{a_i z-a_i^2/2}.
\]
If all $a_i>0$, $S$ is strictly increasing from zero to infinity.
It crosses $k$ exactly once. The sign of $L_k'$ changes from positive to
negative there, so its critical value is the global maximum. The analogous
all-negative case gives a minimum; its supremum is attained at an endpoint.
If zero rates accompany rates of one sign, the constant part changes the
finite asymptotic endpoint of $S$ and must be included in the crossing test.
All-zero rates give a constant $S$ and a monotone or constant $L_k$.

If rates of both signs occur, $S''>0$ on $\mathbb R$,
$S(z)\to\infty$ at both ends, and $S'$ increases from minus infinity to
plus infinity. Thus $S$ has one global minimum, at the unique zero $z_0$
of $S'$. If $S(z_0)>k$, $L_k'<0$ everywhere and the endpoint values
suffice. Equality gives a stationary horizontal tangent without a local
maximum. If $S(z_0)<k$, strict convexity and the two infinite limits give
exactly two roots $z_-<z_0<z_+$ of $S(z)=k$. The derivative signs of
$L_k$ are negative, positive and negative on the successive intervals.
The lower root is therefore a local minimum, and the only possible
finite global maximum is $L_k(z_+)$. Taking the maximum of that value
and the two limiting endpoint values completes the optimisation. This
argument uses the positive coefficients $A_i$, not nonnegative original
asset covariances.

The implementation must retain certified sign brackets for $S'$ and the
relevant roots of $S-k$. Positive endpoint and negative endpoint enclosures
prove the brackets; continuity and strict monotonicity on the relevant
branch prove existence and uniqueness. A sign-ambiguous midpoint is not
a sign decision. Evaluating $L_k$ on the final outward root interval
encloses its critical value. The maximum of candidate lower endpoints
and the maximum of candidate upper endpoints enclose their maximum.
If the minimum or tangent comparison cannot be decided at the declared
precision, a value at any fixed threshold still supplies a valid lower
bound, but it must not be reported as a certified optimum. The finite
producer explicitly records its supported cases and any such fallback.

\subsection{Capped calls and rigorous monetary decoding}

The elementary identity
$\min\{(A-K)_+,H\}=(A-K)_+-(A-K-H)_+$ holds pointwise.
For valid undiscounted call enclosures
$[\ell_K,u_K]$ and $[\ell_{K+H},u_{K+H}]$, subtraction and
$0\leq F_H\leq H$ give
\[
 \max\{0,\ell_K-u_{K+H}\}\leq\mathbb E F_H(A)
 \leq\min\{H,u_K-\ell_{K+H}\}.
\]
This is a derived call-spread bound. Separate-strike bounds need not retain
all cancellation, and the result is not claimed to be the sharpest capped
bound obtainable from the published approach. The intersection cannot be
empty for valid inputs; an empty numerical intersection signals an error.
Discounting multiplies the final nonnegative interval by the positive
outward enclosure of $e^{-r_f\tau}$. The midpoint of its exact rational
endpoints has worst-case absolute error equal to half their difference.
A cent certificate requires that complete discounted radius to be at most
$1/100$. The producer records an undiscounted rational radius before
serialising an Arb interval, whose radius representation can widen its
endpoints slightly. The independent final check also uses the exact saved
outward endpoints; all admission classifications agree in this cohort. It is deterministic and has no sampling failure probability.
If the undiscounted radius also passes, that stricter condition is reported
separately rather than silently substituted for the discounted target.

\subsection{What the same polynomial information permits classically}

For the exact conditioned scalar function $g_m(t)=P_\lambda(m+\log t)$,
Taylor's integral identity at $t=1$ gives the degree-three price centre
$c_0+c_1+c_2+c_3$ and radius $M_4\mathbb E U^4/24$.
It is a finite identity, not an assumption of convergence of an infinite
lognormal Taylor series. The degree-one and degree-two alternatives have
radii $M_2\mathbb E U^2/2$ and $M_3\mathbb E U^3/6$ respectively;
vanilla convexity makes the degree-one interval one-sided. Lower-order
valid intervals can be retained when a higher-order radius grows.

For $C_0=P_\lambda$, $C_1=P_\lambda'$ and the falling boundary
operators of the main paper, the coefficients are
$ c_k=\mathbb E[U^kC_{k,\lambda}(m)]/k!$.
Write $X_i=e^{d_i+s_i\cdot Y}$ and
$q_\alpha=j!\prod_iw_i^{\alpha_i}/\prod_i\alpha_i!$ for
$|\alpha|=j$. Expanding $U^k=(\sum_iw_iX_i-1)^k$ and completing
squares in the Gaussian density gives, with
$v_\alpha=\sum_i\alpha_i s_i$ and
$\sigma^2=\lambda^2+\|R_G\|^2$,
\[
 c_k=\frac1{k!}\sum_{j=0}^k(-1)^{k-j}{k\choose j}
 \sum_{|\alpha|=j}q_\alpha
 e^{\sum_i\alpha_id_i+\|v_\alpha\|^2/2}
 C_{k,\sigma}(m_G+R_G\cdot v_\alpha).
\]
The Gaussian convolution commutes with these constant-coefficient
operators; bounded falling derivatives justify differentiation under the
integral for $k\geq2$, while the integrable exponential growth of the
call kernel supplies domination for $k=0,1$. The same expansion without
the boundary kernel gives the continuous $U^k$ moments. Signed terms
are retained through the finite sum. The complete Gaussian-tilt proof
and independent polynomial expansion remain in the indexed supporting
classical derivations.

The transport-weighted moments in the grid certificate do not literally
equal the full Gaussian moments needed here. Nevertheless, the same model,
positive basket geometry and shared $M_2,M_3,M_4$ permit the latter to be
constructed from finite exponential sums of degrees at most four.
Gaussian tilting also supplies the price coefficients, including the
nonzero shift $R_G\cdot\sum_i\alpha_i s_i$. Those additional kernel
evaluations and Gaussian moments are paid work. Granting the common model
information does not grant a numerical basket price. The bounded
comparison retains these costs and both failures and successes.
Neither a smaller grid-transfer allowance nor a finite-moment formula
alone establishes superiority over these continuous price certificates.
